\section{Methodology}
\label{sec:methodology}

The methodology has three layers: the fixed family-tuple likelihood model, the
penalized expectation-maximization computation used to fit a candidate, and the
model-search device used to compare candidate tuples.  We state these layers
separately because they have different mathematical status and lead to
different inferential claims.

\subsection{Scope of the methodology}

We first define the statistical object studied in the paper and separate three
regimes that require different inferential arguments.

First, the regular fixed-model regime treats the number of components \(K\) and
the ordered tuple of component families
\[
    \mathfrak m=(d_1,\ldots,d_K)
\]
as fixed.  The true parameter is assumed to be an interior point, no two
components coincide, mixing weights are bounded away from zero, and the Fisher
information is nonsingular after a labeling convention is imposed.  The
identifiability, consistency, asymptotic normality, Louis-information, and EM
monotonicity results below are proved for this regime.

Second, penalized estimation is used as a computational and regularization
device.  Penalization changes the distributional limit of the estimator unless
the tuning parameter is asymptotically negligible at the \(n^{-1/2}\) scale.
For this reason, the likelihood-based standard errors in this paper target the
unpenalized estimator, or a post-selection refit on a fixed selected model, not
the raw lasso estimator.

Third, choosing \(K\), choosing the family tuple, and comparing nested families
such as Poisson inside negative binomial are model-selection and nonregular
problems.  We describe the procedure used in the computations, but the regular
asymptotic theorems do not by themselves justify ordinary Wald inference after
family search, component-number search, or boundary selection.

\subsection{Fixed family-tuple model}

Let \((Y_i,X_i)\), \(i=1,\ldots,n\), be independent observations from the joint
law of \((Y,X)\), where \(Y\in\mathcal Y\subseteq\mathbb R\) and
\(X\in\mathcal X\subseteq\mathbb R^p\).  The first component of \(X\) is an
intercept.  For a fixed model \(\mathfrak m=(d_1,\ldots,d_K)\), assume that all
families in \(\mathfrak m\) are dominated by a common sigma-finite measure
\(\nu\) on \(\mathcal Y\).  This excludes mixtures whose component laws live on
incompatible supports unless a common dominating measure has been explicitly
specified.

For family \(d\), let
\[
    f_d(y;\mu,\gamma)
\]
denote its density or mass function with respect to \(\nu\), where
\(\mu\in\mathcal M_d\) is the systematic location or mean-like parameter and
\(\gamma\in\Gamma_d\) is a finite-dimensional nuisance parameter.  The inverse
link is denoted by
\[
    h_d(\eta)=g_d^{-1}(\eta).
\]
For component \(k\), define
\[
    q_k(y\mid x;\beta_k,\gamma_k)
    =
    f_{d_k}\{y;h_{d_k}(x^\top\beta_k),\gamma_k\}.
\]
The conditional mixture density is
\[
    p_\theta(y\mid x)
    =
    \sum_{k=1}^K \pi_k q_k(y\mid x;\beta_k,\gamma_k),
    \qquad
    \pi_k>0,\quad \sum_{k=1}^K\pi_k=1,
\]
where
\[
    \theta=\{(\pi_k,\beta_k,\gamma_k):k=1,\ldots,K\}\in\Theta_{\mathfrak m}.
\]
The family labels \(d_k\) are part of the model \(\mathfrak m\), not continuous
parameters.

Because finite mixtures are invariant under relabeling, \(\theta\) is
identifiable at best up to permutations that preserve the multiset of
components.  In local asymptotic statements we work in a labeled neighborhood of
the true parameter, for example by imposing a deterministic ordering rule.  All
global identifiability statements are understood modulo label permutations.

\subsection{Observed and penalized likelihoods}

For fixed \(\mathfrak m\), the observed log-likelihood is
\[
    \ell_n(\theta)
    =
    \sum_{i=1}^n
    \log p_\theta(Y_i\mid X_i).
\]
The unpenalized estimator is
\[
    \hat\theta_n
    \in
    \argmax_{\theta\in\Theta_{\mathfrak m}}\ell_n(\theta).
\]

For regularized estimation, let \(\beta_{k,-0}\) denote the non-intercept
coordinates of \(\beta_k\).  We use penalties of the form
\[
    P_n(\theta)
    =
    \sum_{k=1}^K \lambda_{n,k}\rho_k(\beta_{k,-0}),
\]
where \(\rho_k\) may be ridge, lasso, or elastic-net.  The penalized estimator
is
\[
    \hat\theta_{n,\lambda}
    \in
    \argmax_{\theta\in\Theta_{\mathfrak m}}
    \{\ell_n(\theta)-P_n(\theta)\}.
\]
The intercept is not penalized.  This convention follows standard sparse
regression practice \citep{tibshirani1996lasso} and is used throughout.

\subsection{Expectation-maximization representation}

Introduce latent variables \(Z_i\in\{1,\ldots,K\}\) with
\(\Pr(Z_i=k)=\pi_k\).  Given \(Z_i=k\), the conditional density of \(Y_i\) is
\(q_k(Y_i\mid X_i;\beta_k,\gamma_k)\).  The complete-data log-likelihood is
\[
    \ell_{c,n}(\theta)
    =
    \sum_{i=1}^n\sum_{k=1}^K
    \mathbbm 1(Z_i=k)
    \{\log\pi_k+\log q_k(Y_i\mid X_i;\beta_k,\gamma_k)\}.
\]
Given the current parameter \(\theta^{(t)}\), the posterior responsibilities are
\[
    \tau_{ik}^{(t)}
    =
    \frac{
        \pi_k^{(t)}q_k(Y_i\mid X_i;\beta_k^{(t)},\gamma_k^{(t)})
    }{
        \sum_{\ell=1}^K
        \pi_\ell^{(t)}
        q_\ell(Y_i\mid X_i;\beta_\ell^{(t)},\gamma_\ell^{(t)})
    }.
\]
The penalized EM objective is
\[
    Q_{P,n}(\theta\mid\theta^{(t)})
    =
    \sum_{i=1}^n\sum_{k=1}^K
    \tau_{ik}^{(t)}
    \{\log\pi_k+\log q_k(Y_i\mid X_i;\beta_k,\gamma_k)\}
    -
    P_n(\theta).
\]
The resulting construction follows the EM principle of
\citet{dempster1977em}.  The mixing proportions have the closed-form update
\[
    \pi_k^{(t+1)}
    =
    n^{-1}\sum_{i=1}^n\tau_{ik}^{(t)}.
\]
The remaining update decomposes by component:
\[
    (\beta_k^{(t+1)},\gamma_k^{(t+1)})
    \in
    \argmax_{\beta_k,\gamma_k}
    \left[
    \sum_{i=1}^n\tau_{ik}^{(t)}
    \log q_k(Y_i\mid X_i;\beta_k,\gamma_k)
    -
    \lambda_{n,k}\rho_k(\beta_{k,-0})
    \right].
\]
When this maximization is exact, the algorithm is an EM algorithm for the
penalized objective.  When the component update is only improved, but not
globally maximized, the algorithm is a generalized EM algorithm.  The
monotonicity result below requires the explicit ascent condition
\[
    Q_{P,n}(\theta^{(t+1)}\mid\theta^{(t)})
    \ge
    Q_{P,n}(\theta^{(t)}\mid\theta^{(t)}).
\]

\subsection{Model search}

Let \(\mathcal D\) be a finite collection of candidate response families.  We
consider only family tuples whose supports are compatible with the observed
support of \(Y\).  For a fitted model \(\mathfrak m\), define
\[
    \mathrm{BIC}(\mathfrak m)
    =
    -2\ell_n(\hat\theta_{\mathfrak m})
    +
    a_{\mathfrak m}\log n,
\]
where BIC denotes the Bayesian information criterion and
\(a_{\mathfrak m}\) is the number of free parameters in the fixed model
\citep{schwarz1978dimensions}.
Penalized fits are used as a screening and regularization device.  For a
lasso-selected or elastic-net-selected support, the final likelihood score is
computed by refitting the same family tuple without penalty while holding the
component-specific active sets fixed.  If
\(\hat S_k=\{j>0:\hat\beta_{kj}\neq0\}\) is the selected non-intercept support
for component \(k\), the refitted model uses the parameter count
\[
    \hat a_{\mathfrak m}
    =
    (K-1)
    +K
    +
    \sum_{k=1}^K
    |\hat S_k|
    +
    \sum_{k=1}^K\dim(\gamma_k),
\]
where the second term counts the \(K\) component intercepts.  We report the
penalized objective and the unpenalized likelihood of the screening fit, but
rank lasso-selected candidates by the BIC of the unpenalized active-set refit
when the refit converges.  Ridge penalties do not define a sparse active set
and are therefore treated as sensitivity analyses rather than as sources of
post-lasso BIC scores.  The refit BIC is still conditional on the selected
support and family tuple; it is a transparent computational score, not a claim
of valid post-selection inference.

For small candidate spaces we enumerate all unordered family tuples.  For
larger spaces the implementation uses a beam search: tuples of size \(K\) are
expanded from retained tuples of size \(K-1\), fitted by the EM procedure, and
only the best few according to the chosen criterion are carried forward.  The
beam search is a deterministic heuristic conditional on initialization and
random seed; it does not inherit the fixed-model regularity theory.

\section{Theoretical properties}
\label{sec:theory}

We next record the likelihood properties that are used to interpret a fitted
candidate model.  The results are intentionally conditional on a fixed,
regular, identifiable candidate; later family search, component-number
selection, and active-set selection are treated as separate nonregular
operations.

\subsection{Regular fixed-model assumptions}

All results in Sections~\ref{sec:identifiability}--\ref{sec:louis_theory}
condition on a fixed family tuple \(\mathfrak m=(d_1,\ldots,d_K)\).  Write
\(\theta_0\) for the true parameter.

\begin{assumption}[Sampling]
\label{ass:sampling}
The observations \((Y_i,X_i)\), \(i=1,\ldots,n\), are independent and
identically distributed.  Conditional on \(X=x\), the true density of \(Y\) is
\(p_{\theta_0}(y\mid x)\) for some \(\theta_0\in\Theta_{\mathfrak m}\).
\end{assumption}

\begin{assumption}[Parameter space]
\label{ass:param_space}
The parameter space \(\Theta_{\mathfrak m}\) is compact after imposing a
labeling convention.  The true parameter \(\theta_0\) is an interior point of
this labeled parameter space.  There exists \(\pi_{\min}>0\) such that
\(\pi_k\ge \pi_{\min}\) for all \(k\) and all \(\theta\in\Theta_{\mathfrak m}\).
Nuisance parameters are bounded away from degeneracies such as zero variances,
infinite dispersions, and boundary zero-inflation probabilities.
\end{assumption}

\begin{assumption}[Continuity and domination]
\label{ass:continuity}
For each \((y,x)\), \(p_\theta(y\mid x)\) is continuous in \(\theta\).  Moreover
\(\log p_\theta(Y\mid X)\) is dominated by an integrable envelope in a
neighborhood of every \(\theta\in\Theta_{\mathfrak m}\), so that
\[
    \sup_{\theta\in\Theta_{\mathfrak m}}
    \left|
    n^{-1}\ell_n(\theta)
    -
    E_{\theta_0}\log p_\theta(Y\mid X)
    \right|
    \to 0
\]
in probability.
\end{assumption}

\begin{assumption}[Primitive component identifiability]
\label{ass:primitive_component}
For each family \(d\), if
\[
    f_d(y;\mu_1,\gamma_1)=f_d(y;\mu_2,\gamma_2)
    \quad\text{for \(\nu\)-almost every }y,
\]
then \((\mu_1,\gamma_1)=(\mu_2,\gamma_2)\).  The inverse link \(h_d\) is
one-to-one on its domain.  Finally, if \(X^\top b=0\) almost surely, then
\(b=0\).
\end{assumption}

\begin{assumption}[Finite-mixture separability]
\label{ass:mixture_separability}
The collection of component regression densities
\[
    \mathcal Q_{\mathfrak m}
    =
    \left\{
    q_k(\cdot\mid\cdot;\beta,\gamma):
    k=1,\ldots,K,\;(\beta,\gamma)\in\mathcal B_{d_k}\times\Gamma_{d_k}
    \right\}
\]
is \(K\)-identifiable in the following sense.  If two mixtures with at most
\(K\) nonzero components satisfy
\[
    \sum_{r=1}^{m} a_r q_r(y\mid x)
    =
    \sum_{s=1}^{m'} b_s \tilde q_s(y\mid x)
\]
for \(P_X\)-almost every \(x\) and \(\nu\)-almost every \(y\), where
\(a_r,b_s>0\), \(\sum_r a_r=\sum_s b_s=1\), and all atoms on each side are
distinct, then \(m=m'\) and the two sets of weighted atoms are equal up to
permutation.
\end{assumption}

\begin{remark}
Assumption~\ref{ass:mixture_separability} is the nontrivial finite-mixture
identifiability condition.  It is intentionally stated explicitly rather than
hidden behind the phrase ``standard mixture arguments.''  It must be verified
or cited for any family tuple used as a formal model.  It can fail at nested or
boundary cases, for example when an overdispersed family contains a simpler
family as a limiting special case; see \citet{teicher1963identifiability} and
\citet{hennig2000identifiability} for foundational results and regression-mixture
qualifications.
\end{remark}

\subsection{Identifiability}
\label{sec:identifiability}

\begin{lemma}[Component regression identifiability]
\label{lem:component_ident}
Under Assumption~\ref{ass:primitive_component}, for a fixed family \(d\),
\[
    f_d\{y;h_d(x^\top\beta_1),\gamma_1\}
    =
    f_d\{y;h_d(x^\top\beta_2),\gamma_2\}
\]
for \(P_X\)-almost every \(x\) and \(\nu\)-almost every \(y\) implies
\[
    \beta_1=\beta_2,
    \qquad
    \gamma_1=\gamma_2.
\]
\end{lemma}

\begin{proof}
By primitive identifiability of the family, for \(P_X\)-almost every \(x\),
\[
    h_d(x^\top\beta_1)=h_d(x^\top\beta_2),
    \qquad
    \gamma_1=\gamma_2.
\]
Since \(h_d\) is one-to-one,
\[
    x^\top(\beta_1-\beta_2)=0
\]
for \(P_X\)-almost every \(x\).  The full-rank condition in
Assumption~\ref{ass:primitive_component} then gives \(\beta_1=\beta_2\).
\end{proof}

\begin{theorem}[Identifiability of the fixed family-tuple model]
\label{thm:fixed_ident}
Suppose Assumptions~\ref{ass:param_space}, \ref{ass:primitive_component}, and
\ref{ass:mixture_separability} hold.  If
\[
    p_{\theta_1}(y\mid x)=p_{\theta_2}(y\mid x)
\]
for \(P_X\)-almost every \(x\) and \(\nu\)-almost every \(y\), then
\(\theta_1\) and \(\theta_2\) are equal up to a permutation of component
labels.
\end{theorem}

\begin{proof}
Write both sides as finite mixtures of distinct component regression densities.
If any repeated atom occurs within either representation, combine its weights;
Assumption~\ref{ass:param_space} keeps all remaining weights positive.
Assumption~\ref{ass:mixture_separability} implies equality of the weighted
atoms up to permutation.  Therefore, after relabeling,
\[
    \pi_k^{(1)}=\pi_k^{(2)}
    \quad\text{and}\quad
    q_k(\cdot\mid\cdot;\beta_k^{(1)},\gamma_k^{(1)})
    =
    q_k(\cdot\mid\cdot;\beta_k^{(2)},\gamma_k^{(2)}).
\]
Lemma~\ref{lem:component_ident} then gives
\((\beta_k^{(1)},\gamma_k^{(1)})
=
(\beta_k^{(2)},\gamma_k^{(2)})\) for every \(k\).
\end{proof}

\subsection{Existence and consistency}

\begin{theorem}[Consistency of the maximum likelihood estimator]
\label{thm:consistency}
Under Assumptions~\ref{ass:sampling}--\ref{ass:mixture_separability}, an
unpenalized maximizer \(\hat\theta_n\) exists with probability tending to one.
Moreover, after applying the labeling convention,
\[
    \hat\theta_n\to\theta_0
\]
in probability.
\end{theorem}

\begin{proof}
Existence follows from compactness of \(\Theta_{\mathfrak m}\) and continuity of
\(\ell_n(\theta)\).  Define the population criterion
\[
    M(\theta)=E_{\theta_0}\log p_\theta(Y\mid X).
\]
By Assumption~\ref{ass:continuity},
\[
    \sup_{\theta\in\Theta_{\mathfrak m}}
    |n^{-1}\ell_n(\theta)-M(\theta)|
    \to 0
\]
in probability.  For any \(\theta\),
\[
    M(\theta)-M(\theta_0)
    =
    -E_{\theta_0}
    \log\frac{p_{\theta_0}(Y\mid X)}{p_\theta(Y\mid X)}
    \le 0,
\]
with equality if and only if
\[
    p_\theta(Y\mid X)=p_{\theta_0}(Y\mid X)
\]
almost surely.  By Theorem~\ref{thm:fixed_ident} and the imposed labeling
convention, equality holds only at \(\theta=\theta_0\).  Hence \(M\) has a
unique maximizer at \(\theta_0\).  The argmax theorem gives
\(\hat\theta_n\to\theta_0\) in probability.
\end{proof}

\subsection{Regular asymptotic normality}

\begin{assumption}[Local differentiability and information]
\label{ass:local_regular}
In a neighborhood of \(\theta_0\), the labeled parameter can be represented as a
vector in an open subset of \(\mathbb R^r\).  The map
\(\theta\mapsto \log p_\theta(Y\mid X)\) is twice continuously differentiable
in this neighborhood, differentiation may be interchanged with expectation, and
the Fisher information
\[
    I(\theta_0)
    =
    E_{\theta_0}
    \left[
        S_{\theta_0}(Y,X)S_{\theta_0}(Y,X)^\top
    \right],
    \qquad
    S_{\theta}(Y,X)=\nabla_\theta\log p_\theta(Y\mid X),
\]
is finite and nonsingular.
\end{assumption}

\begin{theorem}[Asymptotic normality in the regular fixed model]
\label{thm:asymptotic_normality}
Under Assumptions~\ref{ass:sampling}--\ref{ass:local_regular},
\[
    \sqrt n(\hat\theta_n-\theta_0)
    \Rightarrow
    N\{0,I(\theta_0)^{-1}\}.
\]
\end{theorem}

\begin{proof}
The score equation satisfies
\[
    0=\nabla_\theta\ell_n(\hat\theta_n)
\]
with probability tending to one for an interior local maximizer.  Taylor
expansion around \(\theta_0\) gives
\[
    0
    =
    \nabla_\theta\ell_n(\theta_0)
    +
    \nabla_\theta^2\ell_n(\tilde\theta_n)
    (\hat\theta_n-\theta_0),
\]
where \(\tilde\theta_n\) lies between \(\hat\theta_n\) and \(\theta_0\).  By the
central limit theorem,
\[
    n^{-1/2}\nabla_\theta\ell_n(\theta_0)
    \Rightarrow
    N\{0,I(\theta_0)\}.
\]
By consistency and the uniform law of large numbers for the Hessian,
\[
    -n^{-1}\nabla_\theta^2\ell_n(\tilde\theta_n)
    \to
    I(\theta_0)
\]
in probability.  Solving the Taylor expansion and applying Slutsky's theorem
proves the result.
\end{proof}

\subsection{Louis identity and standard errors}
\label{sec:louis_theory}

The preceding asymptotic normality result requires the observed information
matrix for a fitted mixture.  Direct numerical differentiation of the observed
mixture likelihood is possible, but it is unstable when components have
different response families or nuisance parameters on different scales.  We
therefore use \citet{louis1982information}'s identity to express the observed information through
complete-data component scores and Hessians.

For one observation, let \(Z_i\in\{1,\ldots,K\}\) be the latent component label
and write
\[
    c_{ik}(\theta)
    =
    \log \pi_k+\log f_{d_k}\{Y_i;h_{d_k}(X_i^\top\beta_k),\gamma_k\}.
\]
The posterior responsibility is
\[
    \tau_{ik}(\theta)
    =
    \Pr_\theta(Z_i=k\mid Y_i,X_i)
    =
    \frac{\exp\{c_{ik}(\theta)\}}
    {\sum_{\ell=1}^K\exp\{c_{i\ell}(\theta)\}}.
\]
Let
\[
    s_{ik}(\theta)=\nabla_\theta c_{ik}(\theta),
    \qquad
    H_{ik}(\theta)=\nabla_\theta^2 c_{ik}(\theta),
\]
where entries outside component \(k\) are zero except for the
mixing-proportion coordinates.

\begin{theorem}[Louis observed-information identity for a finite mixture]
\label{thm:louis}
Under the regular fixed-model conditions, the observed information contribution
of observation \(i\) is
\[
    I_i(\theta)
    =
    -\sum_{k=1}^K\tau_{ik}(\theta)H_{ik}(\theta)
    -
    \sum_{k=1}^K\tau_{ik}(\theta)
    \{s_{ik}(\theta)-\bar s_i(\theta)\}
    \{s_{ik}(\theta)-\bar s_i(\theta)\}^\top,
\]
where
\[
    \bar s_i(\theta)=\sum_{k=1}^K\tau_{ik}(\theta)s_{ik}(\theta).
\]
Equivalently,
\[
    I_{\mathrm{obs},n}(\theta)
    =
    \sum_{i=1}^n I_i(\theta)
    =
    -E_\theta\{\nabla_\theta^2\ell_c(\theta)\mid Y,X\}
    -
    \operatorname{Var}_\theta\{\nabla_\theta\ell_c(\theta)\mid Y,X\}.
\]
\end{theorem}

\begin{proof}
For fixed \((Y_i,X_i)\), the observed log-likelihood contribution is
\[
    \ell_i(\theta)=\log\sum_{k=1}^K\exp\{c_{ik}(\theta)\}.
\]
Differentiating once gives Fisher's identity in finite-sum form,
\[
    \nabla_\theta\ell_i(\theta)
    =
    \sum_{k=1}^K\tau_{ik}(\theta)s_{ik}(\theta).
\]
Since
\[
    \nabla_\theta\tau_{ik}(\theta)
    =
    \tau_{ik}(\theta)\{s_{ik}(\theta)-\bar s_i(\theta)\},
\]
a second differentiation yields
\[
    \nabla_\theta^2\ell_i(\theta)
    =
    \sum_{k=1}^K\tau_{ik}(\theta)H_{ik}(\theta)
    +
    \sum_{k=1}^K\tau_{ik}(\theta)
    \{s_{ik}(\theta)-\bar s_i(\theta)\}
    \{s_{ik}(\theta)-\bar s_i(\theta)\}^\top.
\]
The observed information is \(-\nabla_\theta^2\ell_i(\theta)\).  Summing over
observations gives the sample identity.
\end{proof}

For the mixing proportions, the implementation uses \(K-1\) free logits
\(a=(a_1,\ldots,a_{K-1})\), with \(a_K=0\) and
\(\pi_k=\exp(a_k)/\sum_{\ell=1}^K\exp(a_\ell)\).  Then, for
\(j=1,\ldots,K-1\),
\[
    \frac{\partial\log\pi_k}{\partial a_j}
    =
    \mathbbm 1(k=j)-\pi_j,
    \qquad
    \frac{\partial^2\log\pi_k}{\partial a\,\partial a^\top}
    =
    -\{\operatorname{diag}(\pi_{1:K-1})
    -\pi_{1:K-1}\pi_{1:K-1}^\top\}.
\]

The component-specific part of \(s_{ik}\) and \(H_{ik}\) is obtained by
combining a family-level derivative block with the GLM link.  For a single
component, write \(\eta=x^\top\beta\), \(\mu=h(\eta)\), and
\(\ell=\log f(y;\mu,\gamma)\).  Define
\[
    a_\mu=\frac{\partial \ell}{\partial\mu},\qquad
    b_{\mu\mu}=\frac{\partial^2\ell}{\partial\mu^2},\qquad
    a_\gamma=\frac{\partial\ell}{\partial\gamma},\qquad
    B_{\gamma\gamma}=\frac{\partial^2\ell}{\partial\gamma\,\partial\gamma^\top},
    \qquad
    b_{\mu\gamma}=\frac{\partial^2\ell}{\partial\mu\,\partial\gamma}.
\]

\begin{proposition}[Regression-coordinate derivative block]
\label{prop:louis_chain_rule}
Let \(\dot h(\eta)\) and \(\ddot h(\eta)\) denote the first and second
derivatives of the inverse link.  The complete-data score and Hessian block in
\((\beta,\gamma)\) coordinates is
\[
    s_\beta=x\,a_\mu\dot h(\eta),
    \qquad
    s_\gamma=a_\gamma,
\]
\[
    H_{\beta\beta}
    =
    xx^\top\{b_{\mu\mu}\dot h(\eta)^2+a_\mu\ddot h(\eta)\},
    \qquad
    H_{\beta\gamma}
    =
    x\{\dot h(\eta)b_{\mu\gamma}\}^\top,
    \qquad
    H_{\gamma\gamma}=B_{\gamma\gamma}.
\]
\end{proposition}

\begin{proof}
The identities follow from the chain rule applied to
\(\ell\{h(x^\top\beta),\gamma\}\).  Since \(x^\top\beta\) is linear in
\(\beta\), the only second derivative introduced by the link is
\(\ddot h(\eta)\) in the scalar \(\eta\)-direction.  Differentiating the
\(\eta\)-score with respect to \(\gamma\) gives the displayed cross block.
\end{proof}

The family-specific quantities
\((a_\mu,b_{\mu\mu},a_\gamma,B_{\gamma\gamma},b_{\mu\gamma})\) are available in
closed form for the families used in the simulations and applications.
Appendix~\ref{app:louis_mstep} lists the analytical blocks and the weighted
M-step equations, while the Supplementary Material gives additional derivation
details and implementation validation.

In the regular unpenalized model, standard errors are obtained from
\[
    \widehat{\operatorname{Var}}(\hat\theta_n)
    =
    I_{\mathrm{obs},n}(\hat\theta_n)^{-1},
\]
where \(I_{\mathrm{obs},n}\) is the sample observed information.  These standard
errors are conditional on the fixed family tuple, fixed number of components,
and fixed labeling convention.

\subsection{Penalization and the limit distribution}

Naive Wald intervals for the raw lasso estimator require care
\citep{knight2000asymptotics,leeb2005selection}.  The next proposition records
the fixed-dimensional limit under an \(\ell_1\) penalty.

\begin{proposition}[Local limit of penalized estimators]
\label{prop:penalty_limit}
Assume the regular fixed-model conditions and let
\(\lambda_n=\max_k\lambda_{n,k}\).  For simplicity write the parameter vector as
\(\theta\) and consider a lasso penalty on selected coordinates.  If
\(\lambda_n/\sqrt n\to\lambda\in[0,\infty)\), then
\[
    \sqrt n(\hat\theta_{n,\lambda}-\theta_0)
\]
converges to the maximizer of
\[
    u^\top Z-\frac12 u^\top I(\theta_0)u
    -
    \lambda
    \sum_{j\in\mathcal P}
    \left[
        u_j\operatorname{sign}(\theta_{0j})\mathbbm 1(\theta_{0j}\neq 0)
        +
        |u_j|\mathbbm 1(\theta_{0j}=0)
    \right],
\]
where \(Z\sim N\{0,I(\theta_0)\}\) and \(\mathcal P\) is the set of penalized
coordinates.  In particular, if \(\lambda=0\) the usual normal limit is
recovered, while if \(\lambda>0\) and any penalized true coefficient is zero,
the limit has a kink and is generally not normal.
\end{proposition}

\begin{proof}
Set \(\theta=\theta_0+u/\sqrt n\).  The local asymptotic quadratic expansion of
the log-likelihood gives
\[
    \ell_n(\theta_0+u/\sqrt n)-\ell_n(\theta_0)
    =
    u^\top n^{-1/2}\nabla\ell_n(\theta_0)
    -
    \frac12 u^\top I(\theta_0)u
    +
    o_p(1)
\]
uniformly on compact sets of \(u\).  The score term converges to
\(Z\sim N\{0,I(\theta_0)\}\).  For a penalized coordinate,
\[
    \lambda_n\left(
        |\theta_{0j}+u_j/\sqrt n|-|\theta_{0j}|
    \right)
    \to
    \lambda u_j\operatorname{sign}(\theta_{0j})
\]
if \(\theta_{0j}\neq0\), and
\[
    \lambda_n|u_j|/\sqrt n\to \lambda |u_j|
\]
if \(\theta_{0j}=0\).  Combining the likelihood and penalty expansions and
using the argmax continuous mapping theorem yields the stated limit.
\end{proof}

\begin{corollary}
If \(\lambda_n/\sqrt n\to0\), the fixed-dimensional lasso penalty is
asymptotically negligible at the \(n^{-1/2}\) scale.  If
\(\lambda_n/\sqrt n\to\lambda>0\), ordinary Hessian-based standard errors for
the penalized estimator do not generally estimate the distribution in
Proposition~\ref{prop:penalty_limit}.  Valid inference after variable selection
therefore requires additional arguments, such as post-selection refitting,
sample splitting, selective inference, bootstrap methods with label alignment,
or a debiasing construction.
\end{corollary}

\subsection{EM monotonicity and stationary points}

\begin{theorem}[Monotonicity of exact and generalized penalized EM]
\label{thm:em_monotonicity}
Let
\[
    F_n(\theta)=\ell_n(\theta)-P_n(\theta).
\]
If an EM update satisfies
\[
    Q_{P,n}(\theta^{(t+1)}\mid\theta^{(t)})
    \ge
    Q_{P,n}(\theta^{(t)}\mid\theta^{(t)}),
\]
then
\[
    F_n(\theta^{(t+1)})\ge F_n(\theta^{(t)}).
\]
\end{theorem}

\begin{proof}
For any \(\theta\) and current value \(\theta^{(t)}\), write
\[
    \ell_n(\theta)
    =
    Q_n(\theta\mid\theta^{(t)})
    -
    E_{\theta^{(t)}}
    \left[
        \log p_\theta(Z\mid Y,X)
        \mid Y,X
    \right],
\]
where \(Q_n\) is the unpenalized auxiliary function.  Subtracting the same
identity evaluated at \(\theta^{(t)}\) yields
\[
    \begin{aligned}
    F_n(\theta)-F_n(\theta^{(t)})
    &=
    Q_{P,n}(\theta\mid\theta^{(t)})
    -
    Q_{P,n}(\theta^{(t)}\mid\theta^{(t)}) \\
    &\quad+
    \operatorname{KL}\{
    p_{\theta^{(t)}}(Z\mid Y,X),
    p_\theta(Z\mid Y,X)
    \}.
    \end{aligned}
\]
The conditional Kullback--Leibler term is nonnegative.  Therefore any update
that does not decrease \(Q_{P,n}\) cannot decrease \(F_n\).
\end{proof}

\begin{proposition}[Limit points of generalized EM]
\label{prop:gem_stationary}
Suppose \(\Theta_{\mathfrak m}\) is compact, \(F_n\) is continuous, and each
component update is a closed ascent map whose fixed points are precisely the
stationary points of \(F_n\) in the appropriate smooth or subgradient sense.
Then every accumulation point of the generalized EM sequence is a stationary
point of \(F_n\).
\end{proposition}

\begin{proof}
By Theorem~\ref{thm:em_monotonicity}, \(F_n(\theta^{(t)})\) is monotone
nondecreasing.  Compactness and continuity imply that this sequence has finite
limits along convergent subsequences.  Let \(\theta^\star\) be an accumulation
point.  If \(\theta^\star\) were not stationary, the closed ascent property
would give a neighborhood of \(\theta^\star\) in which one full generalized EM
update increases \(F_n\) by at least some positive amount.  A subsequence
entering this neighborhood would then force the monotone objective sequence to
increase beyond its limiting value, a contradiction.  Hence every accumulation
point is stationary.
\end{proof}

\subsection{BIC in the regular fixed-candidate regime}

\begin{proposition}[Regular BIC consistency, excluding singular overfitting]
\label{prop:bic_regular}
Consider a finite collection of candidate models.  Suppose the true model is in
the collection, all candidates satisfy the regular fixed-model assumptions, all
incorrect non-nested candidates have a positive Kullback--Leibler gap from the
truth, and any larger model containing the truth is regular with ordinary
Wilks-type \(O_p(1)\) likelihood gain.  Then BIC selects the smallest true model
with probability tending to one.
\end{proposition}

\begin{proof}
For an incorrect model, the maximized average log-likelihood converges to its
pseudo-true population value, which is smaller than the true model value by a
positive Kullback--Leibler gap.  Hence its BIC exceeds the true model BIC by an
order \(n\) term with probability tending to one.  For a regular overfitted
model containing the truth, the maximized log-likelihood gain over the smallest
true model is \(O_p(1)\), while the BIC penalty difference is
\((a_{\mathrm{large}}-a_{\mathrm{true}})\log n\to\infty\).  Thus regular
overfitted models are rejected with probability tending to one.
\end{proof}

\begin{remark}
Proposition~\ref{prop:bic_regular} does not cover the hardest cases for the
present paper: overfitted mixtures, zero mixing weights, coinciding components,
and nested families such as Poisson as a boundary or limiting case of negative
binomial.  These are singular model-selection problems.  The empirical BIC
search used below should therefore be interpreted as a practical selection
procedure unless a separate singular-learning or problem-specific selection
theory is supplied \citep{keribin2000order,drton2017sbic}.
\end{remark}
